% =====================================================================
%  Virtual Sub-States and the Reversible Layer of Loschmidt's Paradox
%
%  A self-contained paper. Every technical term is defined in the text.
%  No reliance on any unpublished companion manuscript.
% =====================================================================
\documentclass[11pt,a4paper]{article}

\usepackage[utf8]{inputenc}
\usepackage[T1]{fontenc}
\usepackage{lmodern}
\usepackage[a4paper,margin=2.4cm]{geometry}
\usepackage{amsmath,amssymb,amsthm,mathtools}
\usepackage{bm}
\usepackage{booktabs}
\usepackage{array}
\usepackage{enumitem}
\usepackage{microtype}
\usepackage[round,sort&compress]{natbib}
\usepackage{tikz}
\usetikzlibrary{arrows.meta,positioning,calc,decorations.pathmorphing,fit}
\usepackage[colorlinks=true,linkcolor=blue!45!black,
            citecolor=blue!45!black,urlcolor=blue!45!black]{hyperref}

% ---- Theorem environments ----
\newtheoremstyle{thmstyle}{\topsep}{\topsep}{\itshape}{}{\bfseries}{.}{.5em}{}
\theoremstyle{thmstyle}
\newtheorem{theorem}{Theorem}[section]
\newtheorem{lemma}[theorem]{Lemma}
\newtheorem{proposition}[theorem]{Proposition}
\newtheorem{corollary}[theorem]{Corollary}
\theoremstyle{definition}
\newtheorem{definition}[theorem]{Definition}
\newtheorem{principle}[theorem]{Principle}
\newtheorem{example}[theorem]{Example}
\theoremstyle{remark}
\newtheorem{remark}[theorem]{Remark}
\newtheorem{notation}[theorem]{Notation}

% ---- Shorthand ----
\newcommand{\R}{\mathbb{R}}
\newcommand{\N}{\mathbb{N}}
\newcommand{\Z}{\mathbb{Z}}
\newcommand{\Cx}{\mathbb{C}}
\newcommand{\Zp}{\mathbb{Z}_{\geq 0}}
\newcommand{\kB}{k_{\mathrm{B}}}
\newcommand{\Mic}{\mu}
\newcommand{\Obs}{\mathcal{O}}
\newcommand{\Count}{N}
\newcommand{\Hist}{\mathcal{H}}
\newcommand{\Vis}{\mathcal{V}}
\newcommand{\phase}{\varphi}
\newcommand{\Tinv}{R}
\newcommand{\flow}{\phi}
\newcommand{\Mmat}{\mathcal{M}}
\newcommand{\half}{\tfrac{1}{2}}
\newcommand{\eps}{\varepsilon}
\newcommand{\Smatrix}{S}

\title{\textbf{Virtual Sub-States and the Reversible Layer\\[2pt]
of Loschmidt's Paradox}}

\author{%
Kundai Farai Sachikonye\\
\small Technical University of Munich\\
\small Chair of Proteomics and Bioanalytics\\
\small \texttt{kundai.sachikonye@wzw.tum.de}}

\date{\today}

\begin{document}
\maketitle

% =====================================================================
\begin{abstract}
\noindent
Loschmidt's reversibility objection observes that the microscopic
equations of motion are invariant under time reversal, while macroscopic
processes are not, and asks how the second law can follow from
time-symmetric mechanics. We argue that the objection conflates two
strata of a system's description that are governed by different logics,
and that separating them dissolves the paradox without recourse to
probabilistic improbability, coarse-graining, or special boundary
conditions. The first stratum is \emph{configurational}: it consists of
the instantaneous microstate together with all decompositions of that
microstate that recover it under a fixed composition law. On this
stratum, evolution is genuinely reversible, and there exists a large
space of \emph{virtual sub-states}---intermediate decompositions whose
components may fall outside the physically admissible region while their
composite remains admissible. The second stratum is \emph{enumerative}:
it consists of the monotone count of completed dynamical events the
system has undergone. We prove that this count cannot decrease, because
any operation that would decrement it is itself an event that increments
it; that the configurational reversibility and the enumerative
monotonicity are logically independent; and, centrally, that virtual
sub-states are available only for the backward completion of an
already-fixed endpoint and never for forward construction---an asymmetry
we make precise as a theorem about composition contexts. Velocity
reversal acts on the configurational stratum and leaves the enumerative
stratum untouched; the second law is the monotonicity of the latter and
is therefore immune to the former. We then exhibit two independent
physical realisations of virtual sub-states---the off-shell internal
lines of quantum field theory, where path opacity is the statement that
only the sum over diagrams is observable, and the phase-coherent
sub-components of a Fourier-transform mass-spectrometry transient, where
a measured fraction of the decomposition lies off the physical
manifold---and show that the nuclear spin echo is a controlled
laboratory test of the reversibility assumption whose seventy-year
verdict is an engineering-independent floor. The framework is presented
as mathematics with physical interpretation sequestered to remarks;
every theorem is proved from stated definitions.
\end{abstract}

\noindent\textbf{Keywords:} Loschmidt's paradox; arrow of time;
irreversibility; reversible dynamics; virtual sub-states; off-shell
states; path opacity; spin echo; partition count; second law of
thermodynamics.

\tableofcontents

% =====================================================================
\section{Introduction}
\label{sec:intro}
% =====================================================================

\subsection{The objection and its persistence}
\label{sec:objection}

In 1876 Josef Loschmidt raised an objection to Boltzmann's programme of
deriving the second law of thermodynamics from the mechanics of atoms
\citep{loschmidt1876}. The microscopic equations of motion---in their
canonical form, Hamilton's equations---possess a symmetry: if one takes
any solution and simultaneously reverses the sign of time and the signs
of all momenta, the result is again a solution. The equations do not
distinguish future from past. Yet macroscopic processes are flagrantly
asymmetric in time: heat flows from hot to cold and never the reverse, a
gas released into a vacuum disperses and never spontaneously
re-concentrates. The objection, in its sharpest form, is this. Let a gas
evolve from a low-entropy macrostate at time $t_0$ to a high-entropy
macrostate at time $T$. Reverse every molecular velocity at $T$. The
reversed state is a legitimate solution of the same equations, and it
will retrace the expansion in reverse, the gas re-collecting itself,
entropy decreasing. The laws permit it. Why is it never seen?

Boltzmann's reply was statistical \citep{boltzmann1877}: the reversed
evolution is not impossible, only overwhelmingly improbable, because the
microstates leading to entropy decrease occupy a vanishing fraction of
the high-entropy macrostate. Over the following century this reply was
refined by Gibbs's coarse-graining \citep{gibbs1902}, the
information-theoretic reformulation \citep{jaynes1957}, the
cosmological-boundary account \citep{penrose1989}, and environmental
decoherence \citep{zurek2003}; the conceptual literature is large
\citep{ehrenfest1959,reichenbach1956,price1996,zeh2007,lebowitz1993}.
These accounts differ in emphasis but share a single structural
commitment, which we will call the \emph{reversibility presupposition}:
that time reversal is a well-defined, in-principle-achievable operation
whose macroscopic consequences merely fail to be realised, for contingent
reasons such as improbability, ignorance, openness, or special initial
conditions. Each account then explains why the reversed evolutions,
though available, are not seen.

\subsection{What this paper does differently}
\label{sec:different}

We do not accept the framing and then explain away its consequences.
We argue that Loschmidt's intuition silently identifies two strata of a
system's description that obey different logics, and that the appearance
of paradox is entirely an artefact of treating them as one. The two
strata are these.

\begin{itemize}[leftmargin=1.4em]
\item The \emph{configurational stratum}: the instantaneous microstate
of the system, together with the entire space of ways of decomposing
that microstate into parts that recompose to it under a fixed
composition law. On this stratum the dynamics is genuinely reversible;
nothing in what follows disputes that. We will moreover show that this
stratum is far larger than the microstate alone, because it contains a
rich family of \emph{virtual sub-states}: decompositions whose individual
components need not be physically admissible, provided their composite
is.

\item The \emph{enumerative stratum}: the integer count of completed
dynamical events---closed oscillation cycles, committed interactions,
recorded distinctions---that the system has undergone since some origin.
This count is the operational content of elapsed time, and it cannot
decrease.
\end{itemize}

The thesis of the paper is that velocity reversal is an operation on the
configurational stratum that leaves the enumerative stratum untouched,
that the second law is the monotonicity of the enumerative stratum, and
that the two are logically independent, so the former cannot affect the
latter. The reason velocity reversal feels as though it should reverse
\emph{everything} is the same reason virtual sub-states feel paradoxical:
both rely on the freedom of the configurational stratum, and both are
silent about the count. We make this precise through a structural
asymmetry---virtual sub-states are available for the backward completion
of a fixed endpoint but never for forward construction
(Section~\ref{sec:asymmetry})---which is the configurational shadow of
the enumerative monotonicity.

Two features distinguish the resulting resolution from the standard ones.
First, it is not probabilistic: nothing here depends on the reversed
microstate being rare. Second, virtual sub-states are not a metaphor but
an identifiable structure with two independent physical realisations,
developed in Sections~\ref{sec:qft} and~\ref{sec:ms}, and the central
experimental signature---the nuclear spin echo---is a controlled
laboratory test of the reversibility presupposition itself
(Section~\ref{sec:echo}), with a definite and repeatedly confirmed
verdict.

\subsection{Organisation}
\label{sec:organisation}

Section~\ref{sec:setup} fixes the mathematical setting: bounded
measure-preserving dynamics, the microstate, the macroscopic observable,
and the event count, with the three kept rigorously distinct.
Section~\ref{sec:count} establishes the monotonicity of the count and the
incoherence of its decrement. Section~\ref{sec:reversible} defines the
configurational stratum and proves its reversibility.
Section~\ref{sec:virtual} introduces virtual sub-states and proves their
existence and abundance. Section~\ref{sec:asymmetry} proves the
forward/backward asymmetry that is the keystone of the resolution.
Section~\ref{sec:resolution} assembles the resolution of Loschmidt's
paradox and contrasts it with the standard accounts.
Section~\ref{sec:qft} develops the quantum-field-theoretic realisation of
virtual sub-states; Section~\ref{sec:ms} the mass-spectrometric one;
Section~\ref{sec:echo} the spin-echo realisation and its experimental
hierarchy. Section~\ref{sec:objections} answers objections;
Section~\ref{sec:conclusion} concludes. Physical interpretation is
sequestered to Remarks throughout; the logical spine is carried entirely
by the formal environments.

% =====================================================================
\section{The Mathematical Setting}
\label{sec:setup}
% =====================================================================

\subsection{Bounded measure-preserving dynamics}
\label{sec:bounded}

We work with a dynamical system in the standard sense of ergodic theory
\citep{walters1982,arnold1989}.

\begin{definition}[Dynamical system]
\label{def:system}
A \emph{dynamical system} is a triple $(\Omega,\mu,\flow_t)$, where
$\Omega$ is a set (the \emph{phase space}) carrying a $\sigma$-algebra of
measurable sets, $\mu$ is a measure on $\Omega$, and
$\flow_t\colon\Omega\to\Omega$, $t\in\R$, is a one-parameter group of
measurable maps (the \emph{flow}) with $\flow_0=\mathrm{id}$ and
$\flow_{t+s}=\flow_t\circ\flow_s$. A point $x\in\Omega$ is a
\emph{microstate}.
\end{definition}

\begin{definition}[Bounded, measure-preserving]
\label{def:boundedmp}
The system $(\Omega,\mu,\flow_t)$ is \emph{bounded} if
$0<\mu(\Omega)<\infty$, and \emph{measure-preserving} if
$\mu(\flow_t^{-1}(A))=\mu(A)$ for every measurable $A$ and every $t$.
\end{definition}

The two conditions are the abstract residue of two ubiquitous physical
facts: that a persistent system occupies a bounded region of its phase
space (an electron does not leave the atom, a gas molecule does not leave
the box, a star does not leave the galaxy), and that Hamiltonian flow
preserves phase-space volume \citep{liouville1838}. We take both as
given; the development below uses only their measure-theoretic content.

\begin{theorem}[Recurrence]
\label{thm:recurrence}
Let $(\Omega,\mu,\flow_t)$ be bounded and measure-preserving. For every
measurable $A\subseteq\Omega$ with $\mu(A)>0$ and every $T>0$, there
exists an integer $n\geq 1$ with $\mu(A\cap\flow_{nT}^{-1}(A))>0$.
Consequently $\mu$-almost every point of $A$ returns to $A$ under
iteration of $\flow_T$.
\end{theorem}

\begin{proof}
This is the Poincar\'e recurrence theorem \citep{poincare1890}; we recall
the proof because the structure recurs below. Fix $T>0$ and write
$f:=\flow_T$. The sets $A,f^{-1}(A),f^{-2}(A),\dots$ all have measure
$\mu(A)$ by measure preservation. If the conclusion failed, then
$\mu(A\cap f^{-n}(A))=0$ for all $n\geq 1$, and an elementary computation
shows the sets $\{f^{-k}(A)\}_{k\geq 0}$ are pairwise disjoint modulo
null sets. Then
$\mu(\Omega)\geq\sum_{k\geq 0}\mu(f^{-k}(A))=\sum_{k\geq 0}\mu(A)=\infty$,
contradicting boundedness. Hence some $n\geq 1$ has
$\mu(A\cap f^{-n}(A))>0$. The pointwise statement follows by the standard
argument \citep{walters1982}.
\end{proof}

\begin{remark}
\label{rem:recurrence}
Recurrence is the first appearance of the tension Loschmidt's objection
exploits. A bounded measure-preserving system returns arbitrarily close
to any prior microstate, infinitely often. If returning the microstate
were the same as returning the system to its past, the second law would
be violated at every recurrence. We will see (Theorem~\ref{thm:distinct})
that it is not.
\end{remark}

\subsection{Time reversal of the flow}
\label{sec:treversal}

\begin{definition}[Time-reversal involution]
\label{def:treversal}
A measurable involution $\Tinv\colon\Omega\to\Omega$ (so
$\Tinv\circ\Tinv=\mathrm{id}$) is a \emph{time reversal} for the flow if
\begin{equation}
\flow_{-t}=\Tinv\circ\flow_t\circ\Tinv
\qquad\text{for all }t\in\R.
\label{eq:treversal}
\end{equation}
For a Hamiltonian system with momenta $p$ and positions $q$, the standard
choice is $\Tinv(q,p)=(q,-p)$, and \eqref{eq:treversal} holds whenever the
Hamiltonian is even in $p$.
\end{definition}

\begin{remark}[The hinge]
\label{rem:hinge}
Equation~\eqref{eq:treversal} is an identity among \emph{solutions} of
the equations of motion: it says that a certain transformation carries one
solution to another. It is silent about whether there exists a
\emph{physical operation}, performable on a system in a laboratory, whose
execution carries the system from the first solution to the second. The
whole of Loschmidt's paradox turns on the slide from the former to the
latter, and the whole of our resolution is the refusal to make that
slide. The configurational stratum is where \eqref{eq:treversal} lives;
the enumerative stratum is where physical operations live; they are
different strata.
\end{remark}

\subsection{Three quantities, kept distinct}
\label{sec:three}

The conceptual core of the setting is a separation that Loschmidt's
intuition collapses. For a system of many constituents we distinguish
three quantities.

\begin{definition}[Microstate, observable, count]
\label{def:three}
Let $(\Omega,\mu,\flow_t)$ be a bounded measure-preserving system.
\begin{enumerate}[label=(\roman*),leftmargin=2em]
\item The \emph{microstate} $\Mic(t):=\flow_t(x_0)\in\Omega$ is the
complete instantaneous specification of the system started from $x_0$.
\item An \emph{observable} is a measurable map $\Obs\colon\Omega\to E$
into some value space $E$; the \emph{observable value} is
$\Obs(\Mic(t))$. An observable is \emph{coarse} if it is many-to-one on a
set of positive measure, so that distinct microstates share an observable
value.
\item The \emph{event count} $\Count(t)\in\Zp$ is the number of completed
dynamical events the system has undergone in $[0,t]$, where an
\emph{event} is the completion of a distinguished, recurrent dynamical
cycle (made precise in Definition~\ref{def:event}).
\end{enumerate}
\end{definition}

\begin{definition}[Event and event count]
\label{def:event}
Fix a measurable \emph{Poincar\'e section} $\Sigma\subset\Omega$ with
$\mu(\Sigma)>0$ that the flow crosses transversally and recurrently
(its existence is guaranteed by Theorem~\ref{thm:recurrence} for any
positive-measure $\Sigma$). An \emph{event} at time $t$ is a transversal
crossing of $\Sigma$ by the trajectory $\Mic(\cdot)$ at $t$. The
\emph{event count} is
\begin{equation}
\Count(t):=\#\{\,s\in(0,t] : \Mic(s)\in\Sigma\text{ is a transversal
crossing}\,\}.
\label{eq:count}
\end{equation}
\end{definition}

\begin{remark}[Why a count, and why it is time]
\label{rem:countistime}
Definition~\ref{def:event} formalises the operational fact that elapsed
time is always read from a counter: a pendulum counts swings, a quartz
watch counts crystal vibrations, the SI second is a stipulated number of
caesium hyperfine cycles. A ``cycle'' here is a topological event---the
trajectory's return to a section---and as such it presupposes no prior
notion of duration; duration is what one \emph{defines} by promoting one
counter to a standard and stipulating how many of its crossings make a
unit. We use $\Count$ as the primitive temporal quantity and recover the
continuous parameter $t$ as the rescaled limit of $\Count$ for a
high-frequency reference. This is why ``time reversal'' will turn out to
be a statement about the flow's solution set
(Remark~\ref{rem:hinge}) and not about $\Count$.
\end{remark}

\begin{principle}[The conflation]
\label{prin:conflation}
Loschmidt's reversibility intuition rests on the tacit identification
\[
\Obs(t)\ \equiv\ \Mic(t)\ \equiv\ \Count(t):
\]
that recovering the observable ($\Obs(2T)=\Obs(0)$) is the same as
restoring the microstate ($\Mic(2T)=\Mic(0)$), which is the same as
returning to the past ($\Count(2T)=\Count(0)$). The three are distinct,
and conflating them is what makes ``the trajectory retraced'' sound like
``the system went back.''
\end{principle}

We will dismantle the conflation in two steps: by proving the count
cannot return (Section~\ref{sec:count}), and by proving that the
configurational reversibility which \emph{does} hold is confined to the
microstate--observable layer (Section~\ref{sec:reversible}).

% =====================================================================
\section{The Enumerative Stratum: Monotonicity of the Count}
\label{sec:count}
% =====================================================================

\subsection{Strict monotonicity}
\label{sec:monotone}

\begin{theorem}[Strict monotonicity of the count]
\label{thm:monotone}
Let $(\Omega,\mu,\flow_t)$ be a bounded, measure-preserving, non-static
system, and let $\Count$ be its event count with respect to a section
$\Sigma$ of positive measure. Then $\Count$ is non-decreasing, and for
every $t_1<t_2$ separated by at least one event, $\Count(t_2)>\Count(t_1)$.
\end{theorem}

\begin{proof}
By Definition~\ref{def:event}, $\Count(t)$ counts crossings in $(0,t]$,
so $t_1<t_2$ gives $\Count(t_1)\leq\Count(t_2)$ (the counted set grows).
By Theorem~\ref{thm:recurrence} applied to $\Sigma$, transversal
crossings recur for almost every trajectory; non-staticity ensures the
trajectory is not pinned off the section. Hence on any interval
$[t_1,t_2]$ containing a crossing, at least one term is added to the
count, so $\Count(t_2)\geq\Count(t_1)+1>\Count(t_1)$.
\end{proof}

\begin{corollary}[Injectivity of the count on time]
\label{cor:inject}
For $t_1\neq t_2$ separated by at least one event, $\Count(t_1)\neq
\Count(t_2)$.
\end{corollary}

\subsection{The incoherence of decrement}
\label{sec:incoherence}

Monotonicity says the count happens not to decrease. The next result says
something stronger and more decisive for the paradox: \emph{no operation
can decrease it}, because the execution of any operation is itself a
counted process.

\begin{theorem}[Incoherence of decrement]
\label{thm:incoherence}
Let $(\Omega,\mu,\flow_t)$ be bounded, measure-preserving, and non-static.
There is no physical operation $\Theta$---understood as a process
occupying a positive interval of the flow---whose execution leaves the
system with $\Count$ strictly smaller than before $\Theta$ began.
\end{theorem}

\begin{proof}
Suppose $\Theta$ begins at time $a$ and ends at time $b>a$ (positive
duration is definitional: an operation that occupies no interval of the
flow performs nothing). During $[a,b]$ the trajectory evolves under
$\flow_t$, and by Theorem~\ref{thm:monotone} it accrues at least the
crossings of $\Sigma$ that occur in $[a,b]$, so
$\Count(b)\geq\Count(a)$, with strict inequality whenever $[a,b]$
contains a crossing. The proposed decrement requires $\Count(b)<\Count(a)$,
which contradicts this inequality. The only escape would be a $\Theta$
with $b\leq a$, i.e.\ one occupying no forward interval, which performs no
operation. Hence no decrementing operation exists.
\end{proof}

\begin{corollary}[The reversal map does not act on the count]
\label{cor:Rnotcount}
The time-reversal involution $\Tinv$ of Definition~\ref{def:treversal}
does not decrement $\Count$. Applying $\Tinv$ is a substitution within the
solution set of the equations of motion
(Remark~\ref{rem:hinge}); it is not a process occupying a forward interval
of the flow, and therefore acts on the configurational data $(q,p)$, not
on $\Count$. Any \emph{physical} realisation of velocity reversal---an
apparatus that flips momenta---is itself an operation with positive
duration and so, by Theorem~\ref{thm:incoherence}, increments $\Count$.
\end{corollary}

\begin{remark}[The arrow of time as a counting fact]
\label{rem:arrow}
Theorem~\ref{thm:incoherence} locates the arrow of time not in a
statistical tendency layered over reversible foundations, but in the
structure of counting itself. The count cannot decrease because counting
is what elapsed time \emph{is} (Remark~\ref{rem:countistime}), and any
attempt to ``go back'' is a further act of counting forward. This is a
sharper statement than improbability: the reversed evolution is not rare,
it is enumeratively foreclosed. We emphasise that nothing here forbids the
\emph{configuration} from retracing its path (Section~\ref{sec:reversible}
shows it may); what is foreclosed is the identification of that retracing
with a decrement of the count.
\end{remark}

\subsection{Visited-cell entropy}
\label{sec:entropy}

We now connect the count to entropy, in a way that makes the second law a
corollary of monotonicity rather than of probability. Fix a countable
measurable partition $\mathcal{P}=\{C_1,C_2,\dots\}$ of $\Omega$ into
cells of positive measure---the abstract image of a finite-resolution
measurement, which can resolve only finitely or countably many distinct
outcomes \citep{shannon1948,cover2006}.

\begin{definition}[Visited set and visited-cell entropy]
\label{def:visited}
The \emph{visited set} of the trajectory up to time $t$ is the family of
cells it has entered,
\begin{equation}
\Vis(t):=\{\,C\in\mathcal{P} : \Mic(s)\in C\text{ for some }s\in[0,t]\,\}.
\label{eq:visited}
\end{equation}
The \emph{visited-cell entropy} is $S(t):=\kB\ln|\Vis(t)|$.
\end{definition}

\begin{theorem}[Monotone entropy]
\label{thm:entropy}
For a bounded, measure-preserving, non-static system,
$\Vis(t_1)\subseteq\Vis(t_2)$ for $t_1\leq t_2$, with strict inclusion on
a set of positive measure of times. Consequently $S(t)$ is non-decreasing,
and strictly increasing whenever the trajectory enters a previously
unvisited cell.
\end{theorem}

\begin{proof}
Inclusion is immediate from \eqref{eq:visited}: a cell entered by time
$t_1$ has been entered by any later $t_2$, and the membership predicate is
never retracted---a cell, once entered, has been entered. Strictness: by
Theorem~\ref{thm:recurrence} the trajectory visits arbitrary
neighbourhoods of its starting point, and by non-staticity it is not
confined to a single cell; hence on a positive-measure set of times it
enters cells not in $\Vis(t_1)$. Each such entry enlarges $\Vis$, so
$|\Vis(t_2)|>|\Vis(t_1)|$ and $S(t_2)>S(t_1)$.
\end{proof}

\begin{corollary}[Second law]
\label{cor:secondlaw}
The visited-cell entropy of a bounded, measure-preserving, non-static
isolated system is a non-decreasing function of time.
\end{corollary}

\begin{remark}[Why this is not Boltzmann's argument]
\label{rem:notboltzmann}
The monotonicity of Theorem~\ref{thm:entropy} is not probabilistic. It
does not say that the visited set typically grows, or grows with high
probability; it says the visited set cannot shrink, because set membership
is monotone under union. The content is combinatorial---a record that is
appended to and never erased---not statistical. Boltzmann's
$H$-theorem~\citep{boltzmann1877}, which quantifies the \emph{rate} of
growth under a molecular-chaos assumption, is a statistical refinement of
this combinatorial skeleton, not its foundation. The reason velocity
reversal cannot decrease $S$ is now visible in advance: reversal
rearranges the microstate, but the set of cells \emph{already visited} is
a fact about the past trajectory, and the past trajectory is not undone by
anything done at the present microstate.
\end{remark}

\begin{remark}[Identity as the unvisited complement]
\label{rem:negation}
The complement $\mathcal{P}\setminus\Vis(t)$---the cells the trajectory
has \emph{not} entered---is a second, equivalent encoding of the visited
set, and its measure shrinks monotonically as $\Vis$ grows. There is a
sense in which a trajectory's identity at time $t$ is carried as much by
what it has excluded as by where it is: two trajectories with the same
current cell but different histories are distinguished precisely by their
visited sets, hence by their complements. This ``identity by exclusion''
will reappear in the spin echo (Section~\ref{sec:echo}), where the
irreversible part of the dynamics is exactly the part that has committed
information to the environment and so cannot be read back from the system
alone.
\end{remark}

% =====================================================================
\section{The Configurational Stratum and Its Reversibility}
\label{sec:reversible}
% =====================================================================

We now turn to the stratum on which reversibility genuinely holds, and
define it with enough structure to host virtual sub-states in
Section~\ref{sec:virtual}.

\subsection{Composition laws and decompositions}
\label{sec:composition}

\begin{definition}[Composition law]
\label{def:complaw}
Let $V$ be a set of \emph{values} (e.g.\ a coordinate range, a momentum, a
field amplitude) equipped with a partial operation
$\boxplus\colon V^k\rightharpoonup V$ for each arity $k\geq 1$, the
\emph{composition law}. A \emph{decomposition} of a value $v\in V$ of
arity $k$ is a tuple $(v_1,\dots,v_k)\in V^k$ with
$v_1\boxplus\cdots\boxplus v_k=v$. We write $\mathrm{Dec}_k(v)$ for the set
of arity-$k$ decompositions of $v$.
\end{definition}

The composition law abstracts the way a physical quantity is assembled
from parts: the way a net coordinate is a mean of sub-coordinates, a total
momentum a sum of constituent momenta, a measured amplitude a sum over
contributing channels. The crucial feature, exploited throughout, is that
$\mathrm{Dec}_k(v)$ is in general very large, and that recomposition
discards which decomposition was used.

\begin{definition}[Configurational state and configurational stratum]
\label{def:configstratum}
Fix a value space $V$, a composition law $\boxplus$, and a
\emph{recomposition observable} $\Obs_V\colon V\to V$ that returns the
value a decomposition recomposes to ($\Obs_V(v)=v$ on values; on a
decomposition $(v_1,\dots,v_k)$ it returns $v_1\boxplus\cdots\boxplus
v_k$). The \emph{configurational state} of the system at time $t$ is the
pair $(\Mic(t),\Obs_V)$. The \emph{configurational stratum} is the space
of all decompositions of all values, quotiented by $\Obs_V$: two
decompositions are configurationally equivalent iff they recompose to the
same value.
\end{definition}

\begin{remark}[What lives where]
\label{rem:whatlives}
The microstate $\Mic(t)$ and any coarse observable $\Obs(\Mic(t))$ live on
the configurational stratum: they are functions of the instantaneous
configuration. The count $\Count(t)$ does not: it is a function of the
\emph{history} of crossings (Definition~\ref{def:event}), not of the
present configuration. This is the formal version of the slogan that the
present configuration carries no timestamp---it cannot, since by recurrence
(Theorem~\ref{thm:recurrence}) the same configuration recurs at
unboundedly many different counts.
\end{remark}

\subsection{Configurational reversibility}
\label{sec:configreversible}

\begin{theorem}[Configurational reversibility]
\label{thm:configreversible}
Let $(\Omega,\mu,\flow_t)$ admit a time reversal $\Tinv$
(Definition~\ref{def:treversal}). Then for every microstate $x$ and every
$t$,
\begin{equation}
\Obs\bigl(\Tinv\flow_t\Tinv\flow_t(x)\bigr)=\Obs(x)
\label{eq:configrev}
\end{equation}
for every observable $\Obs$: the configuration is exactly restored by
evolving forward, reversing, evolving forward again, and reversing once
more. In particular, on the configurational stratum the dynamics is
reversible.
\end{theorem}

\begin{proof}
By \eqref{eq:treversal}, $\flow_{-t}=\Tinv\flow_t\Tinv$, so
$\Tinv\flow_t\Tinv\flow_t=\flow_{-t}\flow_t=\flow_0=\mathrm{id}$. Applying
any $\Obs$ to $\mathrm{id}(x)=x$ gives \eqref{eq:configrev}.
\end{proof}

\begin{remark}[The two strata do not interfere]
\label{rem:noninterference}
Theorem~\ref{thm:configreversible} and Theorem~\ref{thm:incoherence} are
both true and not in tension: the first says the configuration returns to
$x$ under the composite map; the second says that performing that
composite map is a process with positive duration that increments
$\Count$. The composite map restores $\Mic$ and $\Obs$ but advances
$\Count$. The system is in the same configuration at a strictly later
count. This is the precise content of ``the trajectory retraced but the
system did not go back,'' and it is the entire resolution in miniature; the
remaining sections enrich the configurational stratum with virtual
sub-states and supply physical realisations.
\end{remark}

\begin{theorem}[Genuine distinctness of the three quantities]
\label{thm:distinct}
The microstate $\Mic$, the coarse observable $\Obs$, and the count
$\Count$ are logically independent in the following sense.
\begin{enumerate}[label=(\roman*),leftmargin=2em]
\item $\Obs$ may return to a past value while $\Mic$ does not: a coarse
observable is many-to-one, so $\Obs(t_2)=\Obs(t_1)$ does not imply
$\Mic(t_2)=\Mic(t_1)$.
\item $\Mic$ may return (arbitrarily close) to a past value while
$\Count$ does not: by recurrence (Theorem~\ref{thm:recurrence}) the
microstate revisits any neighbourhood, but by monotonicity
(Theorem~\ref{thm:monotone}) the count has strictly increased meanwhile.
\item $\Count$ never returns to a past value (Theorem~\ref{thm:incoherence}).
\end{enumerate}
\end{theorem}

\begin{proof}
(i) is the definition of a coarse observable
(Definition~\ref{def:three}). (ii): at a recurrence time $t_2$ with
$\Mic(t_2)$ near $\Mic(t_1)$, Theorem~\ref{thm:monotone} gives
$\Count(t_2)>\Count(t_1)$ because crossings accrued in the interim. (iii)
is Theorem~\ref{thm:incoherence}.
\end{proof}

\begin{remark}[Dissolution of the recurrence objection]
\label{rem:zermelo}
Theorem~\ref{thm:distinct}(ii) settles Zermelo's recurrence objection to
Boltzmann \citep{zermelo1896,boltzmann1877}: a Poincar\'e recurrence is a
return of $\Mic$, not of $\Count$, and it is $\Count$---through the
visited set $\Vis$ of Theorem~\ref{thm:entropy}---that carries entropy.
The microstate may recur a googol times; the visited set only grows.
\end{remark}

% =====================================================================
\section{Virtual Sub-States}
\label{sec:virtual}
% =====================================================================

We now show that the configurational stratum is far larger than the
microstate, because the decomposition space $\mathrm{Dec}_k(v)$ contains
elements whose components are not themselves admissible configurations.
These are the virtual sub-states. They are the structure that makes the
reversibility of the configurational stratum \emph{do work}---in the
backward completion of Section~\ref{sec:asymmetry} and in the physical
realisations of Sections~\ref{sec:qft}--\ref{sec:ms}.

\subsection{Admissibility and virtuality}
\label{sec:admissibility}

\begin{definition}[Admissible region]
\label{def:admissible}
Within the value space $V$, fix a measurable \emph{admissible region}
$A\subseteq V$: the set of values that correspond to a physically
realisable configuration (a coordinate inside the unit cell, a four-momentum
on the mass shell, a partition coordinate inside $[0,1]$). A value in $A$
is \emph{on-shell}; a value in $V\setminus A$ is \emph{off-shell}.
\end{definition}

\begin{definition}[Virtual sub-state]
\label{def:virtual}
Let $v\in A$ be an admissible value and $(v_1,\dots,v_k)\in\mathrm{Dec}_k(v)$
a decomposition. The decomposition is a \emph{virtual sub-state} of $v$ if
at least one component $v_i$ is off-shell, $v_i\in V\setminus A$. It is a
\emph{physical} (on-shell) decomposition if every component is admissible.
\end{definition}

The defining tension is that a virtual sub-state recomposes to an
admissible value through inadmissible parts. The composition law permits
this precisely because $\Obs_V$ records only the composite, discarding the
parts.

\subsection{Existence and abundance}
\label{sec:abundance}

We make the abundance precise for the two composition laws that recur in
the physical realisations: the mean (Section~\ref{sec:ms}) and the sum
under a linear conservation constraint (Section~\ref{sec:qft}). Take
$V=\R^m$ with the arity-$k$ mean
$\boxplus(v_1,\dots,v_k)=\tfrac{1}{k}\sum_i v_i$ and admissible region
$A=[0,1]^m$.

\begin{theorem}[Existence of virtual sub-states]
\label{thm:virtualexist}
Let $v\in (0,1)^m$ and $k\geq 2$. Then $\mathrm{Dec}_k(v)$ contains
virtual sub-states; indeed the set of $(v_1,\dots,v_k)$ with
$\tfrac1k\sum_i v_i=v$ and at least one $v_i\notin[0,1]^m$ has positive
Lebesgue measure within the constraint affine subspace
$\{\,(v_1,\dots,v_k):\tfrac1k\sum_i v_i=v\,\}$.
\end{theorem}

\begin{proof}
The constraint set is an affine subspace of $(\R^m)^k$ of dimension
$m(k-1)$, carrying a natural Lebesgue measure. The on-shell decompositions
are its intersection with the box $([0,1]^m)^k$, a bounded convex set of
finite measure. The full constraint subspace is unbounded (one may send
$v_1\to+\infty$ along any axis while compensating with $v_2\to-\infty$,
holding the mean fixed), hence of infinite measure. The off-shell
decompositions are the complement of a finite-measure set in an
infinite-measure space and therefore have positive---indeed
infinite---measure. Concretely, for $m=1$, $v\in(0,1)$, $k=2$, the family
$(v+c,\,v-c)$ for $c>1-v$ has first component $>1$, hence off-shell, and is
a one-parameter family of virtual sub-states.
\end{proof}

\begin{theorem}[The composite constrains no single part]
\label{thm:miracle}
For the mean law on $V=\R^m$ with admissible region $[0,1]^m$, fix any
admissible composite $v\in[0,1]^m$, any arity $k\geq 2$, and any
prescribed values $w_1,\dots,w_{k-1}\in\R^m$ for the first $k-1$ parts
(however far off-shell). Then there is a unique $k$-th part
$w_k=k\,v-\sum_{i<k}w_i$ making $(w_1,\dots,w_k)$ a decomposition of $v$.
\end{theorem}

\begin{proof}
The mean constraint $\tfrac1k\sum_i w_i=v$ is a single affine equation in
$w_k$ with solution $w_k=k\,v-\sum_{i<k}w_i$, which exists and is unique.
\end{proof}

\begin{remark}[Local infeasibility, global feasibility]
\label{rem:localglobal}
Theorem~\ref{thm:miracle} is the formal content of the slogan that a
globally admissible quantity imposes \emph{no constraint whatever} on any
proper subset of its parts: one may fix all but one part arbitrarily, and
the remaining part absorbs the constraint. The components may be wildly
off-shell---``locally infeasible''---while their composite is on-shell and
``globally feasible.'' This is exactly the freedom that the reversibility
of the configurational stratum supplies and that the count knows nothing
about. It is also, as the next two physical sections show, not a formal
curiosity but the operative mechanism behind off-shell internal lines in
field theory and phase-coherent sub-components in spectrometry.
\end{remark}

\begin{remark}[Unobservability by recomposition]
\label{rem:pathopacity}
Because $\Obs_V$ records only the composite, two distinct decompositions
of the same admissible value are indistinguishable by any measurement of
that value. We will call this \emph{path opacity}: the recomposed quantity
contains no record of which decomposition produced it. Path opacity is
what makes virtual sub-states unobservable in isolation while remaining
indispensable to the composite---the precise statement, in field theory,
that single Feynman diagrams are not observable but their coherent sum is
(Section~\ref{sec:qft}).
\end{remark}

% =====================================================================
\section{The Forward/Backward Asymmetry}
\label{sec:asymmetry}
% =====================================================================

We can now prove the keystone. Virtual sub-states are available for one
purpose and not its mirror image: they may be used to \emph{complete} a
computation toward a fixed admissible endpoint, but not to \emph{construct}
a trajectory forward step by step. This asymmetry is the configurational
shadow of the enumerative monotonicity of Section~\ref{sec:count}, and it
is what blocks any attempt to use the freedom of the configurational
stratum to run the dynamics backward.

\subsection{Forward construction versus backward completion}
\label{sec:forwardbackward}

\begin{definition}[Forward construction]
\label{def:forward}
A \emph{forward construction} produces a sequence of configurations
$v^{(0)},v^{(1)},\dots,v^{(n)}$ in which each $v^{(j+1)}$ is obtained from
$v^{(j)}$ by a step that must itself be an admissible configuration: each
intermediate $v^{(j)}\in A$.
\end{definition}

\begin{definition}[Backward completion]
\label{def:backward}
A \emph{backward completion} is given a fixed admissible endpoint
$v^{(n)}=v^{\ast}\in A$ and produces, for each level $j$, a decomposition
of $v^{\ast}$'s coordinate at that level, with the only requirement that
each decomposition recompose correctly; intermediate components are
unconstrained and may be virtual (off-shell).
\end{definition}

\begin{theorem}[Forward constructions cannot use virtual sub-states]
\label{thm:forwardnovirtual}
In a forward construction, no $v^{(j)}$ may be a virtual sub-state: every
intermediate is on-shell.
\end{theorem}

\begin{proof}
By Definition~\ref{def:forward}, each $v^{(j)}\in A$. A virtual sub-state
has, by Definition~\ref{def:virtual}, an off-shell component
$v_i\in V\setminus A$; were $v^{(j)}$ itself that component it would lie
outside $A$, contradicting admissibility of the intermediate. Hence no
intermediate of a forward construction is virtual.
\end{proof}

\begin{theorem}[Backward completions may use virtual sub-states]
\label{thm:backwardvirtual}
In a backward completion toward an admissible endpoint $v^{\ast}$, the
intermediate decompositions may be virtual; moreover, by
Theorem~\ref{thm:virtualexist} the set of completions whose intermediates
are virtual has positive measure, and by Theorem~\ref{thm:miracle} the
endpoint constrains no single intermediate component.
\end{theorem}

\begin{proof}
Immediate from Definition~\ref{def:backward}: intermediate components are
required only to recompose to $v^{\ast}$, not to be admissible. Existence
and abundance are Theorems~\ref{thm:virtualexist} and~\ref{thm:miracle}.
\end{proof}

\begin{theorem}[Asymmetry]
\label{thm:asymmetry}
The use of virtual sub-states is exclusive to backward completion: a
process that admits virtual intermediates is a backward completion toward a
fixed endpoint and is not a forward construction; equivalently, restricting
a backward completion to on-shell intermediates collapses it to a forward
construction.
\end{theorem}

\begin{proof}
By Theorem~\ref{thm:forwardnovirtual} a forward construction admits no
virtual intermediate, so any process with a virtual intermediate is not a
forward construction. By Theorem~\ref{thm:backwardvirtual} a backward
completion does admit them. Restricting a backward completion to on-shell
intermediates imposes $v^{(j)}\in A$ at every level, which is precisely
Definition~\ref{def:forward}; the resulting object is a forward
construction.
\end{proof}

\begin{remark}[Why the asymmetry is the resolution's keystone]
\label{rem:keystone}
The reason this matters for Loschmidt is the following. The configurational
stratum is reversible (Theorem~\ref{thm:configreversible}), and it is rich,
because of virtual sub-states (Section~\ref{sec:virtual}). One might hope to
exploit this richness to drive the system backward---to ``complete''
the present configuration back into a past one through clever, possibly
off-shell, intermediate decompositions. Theorem~\ref{thm:asymmetry} says
this hope is misdirected: virtual sub-states complete toward a \emph{fixed}
endpoint; they do not run a trajectory in either temporal direction. The
endpoint they complete toward is whatever admissible configuration the
system is \emph{already at}. The freedom is entirely internal to a single
count---it is freedom in how the present configuration may be analysed,
not freedom to change which count one is at. The count advances regardless
(Theorem~\ref{thm:incoherence}). Thus the richest available reversible
structure is powerless against the enumerative arrow: it operates within a
count, never across counts.
\end{remark}

\subsection{Cost of the on-shell restriction}
\label{sec:cost}

The asymmetry has a quantitative face: completions that forgo virtual
sub-states lose the efficiency the virtual ones provide. We record the
combinatorial fact for the ternary case, which recurs physically in
Section~\ref{sec:ms}.

\begin{proposition}[Decomposition multiplicity]
\label{prop:multiplicity}
The number of ordered decompositions (compositions) of a positive integer
$n$ into a sequence of positive parts is $2^{n-1}$; the number of such
compositions whose parts are additionally labelled by one of $d$ channels
is $d\,(d+1)^{n-1}$.
\end{proposition}

\begin{proof}
A composition of $n$ corresponds bijectively to a choice of a subset of the
$n-1$ gaps between $n$ unit dots at which to place a divider, giving
$2^{n-1}$ \citep{stanley2011}. Labelling each of the $k$ parts of a
$k$-part composition by one of $d$ channels multiplies the count for that
$k$ by $d^{k}$; summing over $k$,
$\sum_{k=1}^{n}\binom{n-1}{k-1}d^{k}=d\sum_{j=0}^{n-1}\binom{n-1}{j}d^{j}
=d(1+d)^{n-1}$ by the binomial theorem.
\end{proof}

\begin{remark}
\label{rem:multiplicity}
The exponential multiplicity of Proposition~\ref{prop:multiplicity} is the
size of the decomposition space a backward completion may range over;
restricting to on-shell intermediates prunes it to the (generally much
smaller) admissible subset. In the embedding of Section~\ref{sec:ms} this
is the difference between $\mathcal{O}(\log N)$ and $\mathcal{O}(N)$ work
to resolve a feature at depth $\log_d N$. The same multiplicity underlies
the abundance of off-shell decompositions in
Theorem~\ref{thm:virtualexist}.
\end{remark}

% =====================================================================
\section{The Resolution of Loschmidt's Paradox}
\label{sec:resolution}
% =====================================================================

We assemble the resolution.

\begin{theorem}[Velocity reversal cannot decrease entropy]
\label{thm:resolution}
Let $(\Omega,\mu,\flow_t)$ be bounded, measure-preserving, non-static, and
admit a time reversal $\Tinv$. Let the system evolve from $t_0$ to $T$ and
let an apparatus then perform a physical velocity reversal. Then:
\begin{enumerate}[label=(\roman*),leftmargin=2em]
\item the configuration thereafter retraces, and any coarse observable can
return to a past value (Theorem~\ref{thm:configreversible});
\item the event count strictly increases throughout, including across the
reversal operation (Theorems~\ref{thm:monotone}, \ref{thm:incoherence},
Corollary~\ref{cor:Rnotcount});
\item the visited-cell entropy $S$ does not decrease, because the visited
set is a monotone function of the trajectory's history and the reversal
acts only on the present configuration (Theorem~\ref{thm:entropy},
Remark~\ref{rem:notboltzmann}).
\end{enumerate}
Consequently velocity reversal restores configuration without restoring
the past: the system is in a retraced configuration at a strictly later
count and strictly larger (or equal) entropy. No decrease of total entropy
occurs.
\end{theorem}

\begin{proof}
(i) is Theorem~\ref{thm:configreversible}. (ii): the pre-reversal
evolution increments $\Count$ by Theorem~\ref{thm:monotone}; the reversal
apparatus is a physical operation of positive duration, so by
Theorem~\ref{thm:incoherence} it too increments $\Count$, and the
mathematical involution $\Tinv$ does not act on $\Count$ by
Corollary~\ref{cor:Rnotcount}; the post-reversal evolution increments
$\Count$ again. (iii): by Theorem~\ref{thm:entropy}, $\Vis$ only grows; the
reversal changes $\Mic$ but not the family of cells already entered, so $S$
cannot fall. The final clause combines (i)--(iii).
\end{proof}

\begin{corollary}[Dissolution of the paradox]
\label{cor:dissolution}
Loschmidt's reversal is neither improbable nor forbidden: the configuration
genuinely retraces. What does not happen is the identification of that
retracing with a return to the past, because the count---the operational
content of elapsed time---strictly increases and carries the monotone
visited-set entropy with it. The paradox is the conflation of
Principle~\ref{prin:conflation}; separating the strata dissolves it.
\end{corollary}

\begin{remark}[Comparison with the standard resolutions]
\label{rem:comparison}
Table~\ref{tab:comparison} situates the present account. Each standard
resolution grants the reversibility presupposition---that time reversal is
a coherent, in-principle-achievable operation---and then explains why its
consequences are not seen. The present account denies the presupposition
at its root: the operation is coherent on the configurational stratum
(where it changes nothing about the arrow) and has no realisation on the
enumerative stratum (where the arrow lives). The denial is not an extra
assumption; it is Theorem~\ref{thm:incoherence}, which follows from
boundedness, measure preservation, and the definition of an operation as a
process of positive duration.
\end{remark}

\begin{table}[t]
\centering
\small
\renewcommand{\arraystretch}{1.25}
\begin{tabular}{p{0.20\textwidth}p{0.40\textwidth}p{0.30\textwidth}}
\toprule
\textbf{Account} & \textbf{Why reversal is not seen} &
\textbf{Residual commitment} \\
\midrule
Boltzmann statistical \citep{boltzmann1877} &
Reversed microstates are astronomically rare. &
Reversal possible, merely improbable. \\
Gibbs coarse-graining \citep{gibbs1902} &
Coarse-grained entropy grows; fine-grained is constant. &
Irreversibility is observer-relative. \\
Information-theoretic \citep{jaynes1957} &
Entropy is ignorance of the microstate. &
Irreversibility is epistemic. \\
Cosmological boundary \citep{penrose1989} &
A low-entropy initial condition sets the gradient. &
Explains the global gradient, not local reversal. \\
Decoherence \citep{zurek2003} &
Entanglement with the environment destroys phase relations. &
Closed-system reversal granted as coherent. \\
\midrule
\textbf{This work} &
Configuration retraces; the count and its visited-set entropy do not
decrease, because reversal acts only on the present configuration. &
None: reversal is coherent configurationally and absent enumeratively
(Thm.~\ref{thm:incoherence}). \\
\bottomrule
\end{tabular}
\caption{The standard resolutions accept the reversibility presupposition
and explain its consequences away; the present account separates the
configurational and enumerative strata and locates the arrow in the
latter.}
\label{tab:comparison}
\end{table}

% =====================================================================
\section{Physical Realisation I: Off-Shell States in Quantum Field Theory}
\label{sec:qft}
% =====================================================================

The abstract virtual sub-state of Section~\ref{sec:virtual} has a precise,
long-established physical instance: the off-shell internal line of a
Feynman diagram. We develop the correspondence and show that path opacity
(Remark~\ref{rem:pathopacity}) is exactly the statement that the sum over
diagrams, not the individual diagram, is observable.

\subsection{The off-shell line as a virtual sub-state}
\label{sec:offshell}

In perturbative quantum field theory \citep{peskin1995,weinberg1995,%
itzykson1980}, a scattering amplitude for external particles of momenta
$p_1,\dots,p_n$ is computed as a sum over Feynman diagrams
\citep{feynman1949,dyson1949}. External lines carry on-shell momenta,
$p_i^2=m_i^2$ (in units $\hbar=c=1$); internal lines carry momenta $q$
constrained only by momentum conservation at each vertex, and generically
$q^2\neq m^2$---they are \emph{off-shell}, or virtual.

\begin{proposition}[Internal lines are virtual sub-states]
\label{prop:internalvirtual}
Take the value space $V$ to be four-momentum space $\R^{1,3}$, the
admissible region $A=\{q:q^2=m^2\}$ the mass shell, and the composition law
the sum constrained by conservation at each vertex. Then the internal
lines of any diagram contributing to a fixed on-shell external
configuration form a virtual sub-state of that configuration in the sense
of Definition~\ref{def:virtual}: their momenta recompose (via conservation)
to the on-shell external data, while each internal momentum is generically
off-shell.
\end{proposition}

\begin{proof}
Momentum conservation at the vertices is a system of linear constraints
expressing the external momenta as a signed sum of internal and external
line momenta; this is the composition law $\boxplus$ specialised to the
constrained sum. The external data $\{p_i\}$ lie in $A$ by construction
(asymptotic states are on-shell). For a loop diagram the loop momentum is
integrated over all of $\R^{1,3}$, the overwhelming part of which has
$q^2\neq m^2$; for a tree internal line the momentum is fixed by
conservation and generically off-shell. Hence the internal configuration
is a decomposition of the external data with off-shell components: a
virtual sub-state.
\end{proof}

\begin{remark}[The propagator weights the off-shell excursion]
\label{rem:propagator}
The Feynman propagator $i/(q^2-m^2+i\eps)$ assigns each internal line a
weight that is largest near the mass shell and falls off as the line goes
further off-shell; the $i\eps$ prescription fixes the contour. The
off-shell excursion is thus not forbidden but weighted---exactly as in the
abstract picture, where off-shell decompositions are permitted and
abundant (Theorem~\ref{thm:virtualexist}) but the dynamics supplies a
weighting over them. For an unstable particle the propagator denominator
$q^2-m^2+im\Gamma$ carries an imaginary part $im\Gamma$ encoding the decay
rate; the real part is the off-shell distance and the imaginary part the
rate of committed transition.
\end{remark}

\subsection{Diagrams as backward completions; path opacity}
\label{sec:diagramsbackward}

\begin{proposition}[A diagram is a backward completion]
\label{prop:diagrambackward}
A Feynman diagram contributing to a fixed on-shell external configuration
$v^{\ast}=\{p_i\}$ is a backward completion (Definition~\ref{def:backward})
of $v^{\ast}$: the endpoint (external lines) is fixed and on-shell, the
intermediate components (internal lines) are unconstrained except by
recomposition (conservation), and they may be virtual.
\end{proposition}

\begin{proof}
The external data are fixed and admissible; the internal data satisfy only
the recomposition constraint and are generically off-shell
(Proposition~\ref{prop:internalvirtual}). This is exactly
Definition~\ref{def:backward}.
\end{proof}

\begin{theorem}[Path opacity and the sum over diagrams]
\label{thm:qftopacity}
Two diagrams $\Gamma_1,\Gamma_2$ with the same external lines but different
internal topologies are distinct backward completions of the same
on-shell endpoint. No measurement performed on the external (asymptotic)
states distinguishes them; only their coherent sum
$\Mmat=\sum_{\Gamma}\Mmat_{\Gamma}$ enters the observable
$|\Mmat|^2$.
\end{theorem}

\begin{proof}
By Proposition~\ref{prop:diagrambackward} each $\Gamma_a$ is a backward
completion of the same endpoint $v^{\ast}$; by path opacity
(Remark~\ref{rem:pathopacity}), the recomposed quantity---here the
$S$-matrix element evaluated on the external momenta---records only the
composite, not which internal decomposition produced it. The observable
transition probability is $|\Mmat|^2=|\sum_\Gamma\Mmat_\Gamma|^2$, which
contains interference cross terms $\Mmat_{\Gamma_1}^{\ast}\Mmat_{\Gamma_2}$;
the individual $\Mmat_\Gamma$ cannot be recovered from $|\Mmat|^2$ without
further structure. Hence single diagrams are not observable and their sum
is.
\end{proof}

\begin{remark}[Unitarity as recomposition bookkeeping]
\label{rem:unitarity}
The optical theorem,
$2\,\mathrm{Im}\,\Mmat(p\!\to\!p)=\sum_f\int d\Pi_f\,|\Mmat(p\!\to\!f)|^2$
\citep{peskin1995,cutkosky1960}, is the statement that the imaginary part
of the forward amplitude equals the total transition rate summed over final
states. In the present language it is a recomposition identity: the rate at
which probability flows out of the initial channel through (off-shell)
virtual intermediates equals the rate at which it is committed into
(on-shell) final channels. The off-shell sub-states are the reversible,
path-opaque interior; the committed final-state rate is the
enumerative quantity. Unitarity is the conservation law tying the two,
the field-theoretic counterpart of the strict bookkeeping between the
configurational and enumerative strata.
\end{remark}

\begin{remark}[Why virtual particles are not observed]
\label{rem:notobserved}
Theorem~\ref{thm:qftopacity} explains the textbook statement that virtual
particles cannot be directly observed \citep{peskin1995}: a measurement of
the external particles cannot determine which internal lines were
exchanged, because the observable depends on the recomposed amplitude and
not on the decomposition. This is path opacity, not a limitation of
apparatus. The off-shell line is a genuine, weighted element of the
reversible configurational interior; its unobservability in isolation is
the same fact as the unobservability of a single decomposition of a
recomposed value (Remark~\ref{rem:pathopacity}).
\end{remark}

% =====================================================================
\section{Physical Realisation II: Phase-Coherent Sub-Components in
Fourier Mass Spectrometry}
\label{sec:ms}
% =====================================================================

The second realisation is metrological and, unlike the field-theoretic
one, directly furnishes a measured off-shell fraction. An ion oscillating
in a Fourier-transform mass analyser is a bounded oscillatory system whose
transient encodes, in its phase and amplitude structure, decompositions
that are partly off-shell.

\subsection{The oscillating ion as a bounded system}
\label{sec:ion}

In an Orbitrap analyser \citep{makarov2000} an ion of mass-to-charge ratio
$m/z$ executes harmonic axial motion at angular frequency
$\omega_z=\sqrt{q\kappa/m}$, with $\kappa$ the field curvature; in a
Fourier-transform ion-cyclotron-resonance cell \citep{comisarow1974} the
relevant frequency is the cyclotron frequency $\omega_c=qB/m$. In both, the
ion is a bounded oscillator in the sense of Definition~\ref{def:boundedmp},
the image current it induces is recorded as a time-domain transient, and
the mass is inferred from the oscillation frequency---never measured
directly, but read from the rate at which the oscillation cycles accrue.

\begin{remark}[Frequency as a count rate]
\label{rem:freqcount}
That mass is read from a frequency is the metrological face of
Remark~\ref{rem:countistime}: the analyser is a counter, and the
``measurement'' of $m/z$ is a comparison of the ion's cycle-accrual rate
against the instrument's reference oscillator. The transient is a record of
counts, and the event count $\Count$ of Definition~\ref{def:event} is, for
this system, literally the number of image-current cycles.
\end{remark}

\subsection{Phase coherence and sub-components}
\label{sec:phasecoherence}

Consider a precursor ion and its dissociation products, all oscillating in
the same trap field. Because the products are produced from the precursor's
continuous oscillatory trajectory, their phases are coherent with it.

\begin{proposition}[Frequency relation of products]
\label{prop:productfreq}
In a fixed Orbitrap field (fixed $\kappa$, fixed charge), a product ion of
mass $m_f$ oscillates at
\begin{equation}
\omega_f=\omega_{\mathrm{prec}}\sqrt{m_{\mathrm{prec}}/m_f},
\label{eq:productfreq}
\end{equation}
where $\omega_{\mathrm{prec}}$ and $m_{\mathrm{prec}}$ are the precursor's
frequency and mass.
\end{proposition}

\begin{proof}
From $\omega_z=\sqrt{q\kappa/m}$ with common $q,\kappa$, dividing the
product and precursor relations gives
$\omega_f/\omega_{\mathrm{prec}}=\sqrt{m_{\mathrm{prec}}/m_f}$.
\end{proof}

\begin{remark}[Sub-components in the transient]
\label{rem:transientcomponents}
By linearity of image-current superposition, a transient containing the
precursor and a coherent product carries both a component at
$\omega_{\mathrm{prec}}$ and one at $\omega_f$ given by
\eqref{eq:productfreq}, the latter phase-locked to the former. The
information that conventional processing discards---phase and the
sub-harmonic amplitude structure---is exactly the decomposition data: the
record of how the precursor's oscillatory state decomposes into coherent
sub-components.
\end{remark}

\subsection{The mean-recovery decomposition and its off-shell fraction}
\label{sec:meanrecovery}

We now exhibit the virtual sub-state structure directly. Normalise the
relevant ion parameters (frequency, charge state, polarity, time) onto a
common scale so that the physical ion corresponds to an admissible
coordinate $v\in[0,1]^m$, and consider decompositions $(v_1,\dots,v_k)$
with mean $v$ (Definition~\ref{def:complaw}, mean law).

\begin{proposition}[Off-shell sub-components exist and recompose]
\label{prop:msoffshell}
For a normalised ion coordinate $v\in(0,1)^m$, the decomposition space
contains mean-recovery tuples with off-shell components (some
$v_i\notin[0,1]^m$), by Theorem~\ref{thm:virtualexist}; their mean recovers
the physical coordinate by construction. These off-shell sub-components are
the spectrometric virtual sub-states.
\end{proposition}

\begin{proof}
Direct specialisation of Theorem~\ref{thm:virtualexist} with the mean law
and admissible region $[0,1]^m$.
\end{proof}

\begin{remark}[A measured off-shell fraction]
\label{rem:measuredfraction}
The content of Proposition~\ref{prop:msoffshell} is empirically accessible:
constructing the stacked decomposition of a real ion across its natural
parameters and computing the fraction of sub-components falling outside
$[0,1]^m$ yields a definite, nonzero off-shell fraction, while the mean of
the components remains the physical (admissible) coordinate in every case.
A nonzero such fraction is the direct measurement of virtuality: it shows
that the reversible configurational interior of a single measured ion
contains genuinely off-shell parts whose composite is on-shell. This is the
metrological counterpart of the field-theoretic statement
(Section~\ref{sec:qft}) that internal lines are off-shell while external
lines are on-shell.
\end{remark}

\begin{remark}[Tailored-waveform cancellation]
\label{rem:swift}
Stored-waveform inverse-Fourier-transform excitation
\citep{marshall1985}, which isolates a target ion by ejecting others, is
the operational use of off-shell sub-components: the ejection waveform is
the inverse transform of a decomposition designed so that the unwanted
ions' contributions cancel in recomposition while the target's survives.
That such a waveform exists and works is the engineering proof that
mean-recovery decompositions with cancelling (and individually off-shell)
parts are physically constructible.
\end{remark}

\subsection{The efficient embedding}
\label{sec:embedding}

The ternary refinement underlying the normalised coordinate embeds the
discrete partition of resolutions into a continuous metric space, where
backward completion is efficient.

\begin{proposition}[Fisher embedding and backward efficiency]
\label{prop:fisher}
The hierarchy of normalised coordinates with ternary refinement embeds into
$[0,1]^{3}$ with the Fisher-information metric
$ds^2=\sum_{i}(dS^i)^2/(S^i(1-S^i))$, the unique metric invariant under the
refinement (a consequence of \v{C}encov's uniqueness theorem
\citep{cencov1982,amari2016}). Backward completion of a feature at
refinement depth $\log_3 N$ then costs $\mathcal{O}(\log_3 N)$ steps,
against $\Omega(N)$ for an on-shell-only search.
\end{proposition}

\begin{proof}[Proof sketch]
Assign each depth-$n$ cell its base-3 address along each axis and map the
address to the corresponding ternary fraction; the induced map is injective
with dense image, and the Fisher metric is the unique refinement-invariant
Riemannian metric on the resulting simplex coordinates
\citep{cencov1982,amari2016}. Descending to a leaf at depth $\log_3 N$ is
$\log_3 N$ ternary choices; an on-shell-only search among $N$ leaves is
$\Omega(N)$, recovering the multiplicity gap of
Remark~\ref{rem:multiplicity}.
\end{proof}

\begin{remark}
\label{rem:embeddingasymmetry}
Proposition~\ref{prop:fisher} is the quantitative form of the asymmetry
(Theorem~\ref{thm:asymmetry}) in the spectrometric setting: the
$\mathcal{O}(\log N)$ efficiency is available only to the backward
completion that may pass through off-shell sub-coordinates; restricting to
on-shell intermediates forfeits it and returns the cost to $\Omega(N)$.
\end{remark}

% =====================================================================
\section{The Spin Echo: A Controlled Test of the Reversibility Assumption}
\label{sec:echo}
% =====================================================================

The two realisations above instantiate virtual sub-states; the spin echo
supplies the experimental verdict on the reversibility presupposition
itself. It is the closest laboratory approximation to Loschmidt's
operation, performed routinely since 1950 \citep{hahn1950}, and its
seventy-year refinement exposes an irreducible floor rather than improving
reversal.

\subsection{The echo as configurational recovery}
\label{sec:echorecovery}

In a spin-echo experiment an ensemble of nuclear spins is tipped into the
transverse plane, allowed to dephase under a spread of precession
frequencies, and then subjected to a refocusing ($180^\circ$) pulse that
inverts each accumulated phase $\phase_i\mapsto-\phase_i$; at twice the
dephasing time the ensemble rephases and the macroscopic signal reappears
as an echo \citep{hahn1950}.

\begin{proposition}[The echo recovers configuration, not history]
\label{prop:echorecovery}
For the static (inhomogeneous) part of the dephasing, the refocusing pulse
restores the observable transverse magnetisation
$\Obs(2\tau)=\Obs(0)$ while the microstate is never lost (each phase is a
determined function of the fixed local offset and the pulse history) and
the event count advances monotonically across $[0,2\tau]$.
\end{proposition}

\begin{proof}
A spin with fixed offset $\Delta\omega_i$ accumulates phase
$\Delta\omega_i\tau$ before the pulse; the pulse sends it to
$-\Delta\omega_i\tau$, and it accumulates $\Delta\omega_i\tau$ again in
$[\tau,2\tau]$, returning to zero phase at $2\tau$ independently of
$\Delta\omega_i$. The vector sum---the observable---is thereby restored:
$\Obs(2\tau)=\Obs(0)$. Each phase is at every instant a determined function
of the fixed offset and the (deterministic) pulse history, so the microstate
is never lost; this is configurational reversibility
(Theorem~\ref{thm:configreversible}) realised in the laboratory. The
spectrometer's reference oscillator crosses its section throughout, so by
Theorem~\ref{thm:monotone} the count strictly increases.
\end{proof}

\begin{remark}[Echo $=$ retraced configuration at a later count]
\label{rem:echolater}
Proposition~\ref{prop:echorecovery} is Remark~\ref{rem:noninterference}
made experimental: the signal returns (configuration restored) while the
count advances (no return to the past). ``The signal came back'' is not
``the system went back.'' One can point at all three quantities at once on
the oscilloscope: the observable that refocused, the microstate that ran
forward intact, and the count that advanced and cannot return.
\end{remark}

\subsection{What cannot be refocused, and why}
\label{sec:irreducible}

The echo never fully recovers. The fluctuating (homogeneous) part of the
dephasing, governed by the transverse relaxation time $T_2$
\citep{bpp1948}, is not refocused, and the echo amplitude is reduced by
$e^{-2\tau/T_2}$.

\begin{proposition}[Refocusability excludes committed interaction]
\label{prop:refocus}
A dephasing contribution is refocusable by a $180^\circ$ pulse if and only
if the phase it produces after the pulse is a determined function of the
present microstate. A contribution arising from a fluctuating interaction
that has committed which-spin information to the environment is not such a
function and is not refocusable.
\end{proposition}

\begin{proof}
If the post-pulse phase is determined by the present microstate, the pulse
$\phase\mapsto-\phase$ pre-compensates it and the contribution refocuses
(the static case of Proposition~\ref{prop:echorecovery}). Conversely, a
fluctuating interaction imprints a phase increment in $[\tau,2\tau]$ that is
statistically independent of the increment in $[0,\tau]$ once which-spin
information has passed into the environment; knowing the present microstate
of the spin does not determine it, so no operation on the spin alone can
pre-compensate it, and the residual randomised phase degrades the
realignment.
\end{proof}

\begin{corollary}[The irreducible floor is the committed record]
\label{cor:floor}
The unrefocusable part of the dephasing is exactly the part that has
committed information to the environment---entered new cells of the joint
system--environment partition in the sense of
Definition~\ref{def:visited}---and the echo deficit $e^{-2\tau/T_2}$ is a
calibrated readout of that committed record. It is non-negative, monotone
in $\tau$, and not reduced by any further pulse.
\end{corollary}

\begin{proof}
By Proposition~\ref{prop:refocus} the unrefocusable contributions are the
committed ones; committing information is entering previously unvisited
cells of the joint partition, which by Theorem~\ref{thm:entropy} only
accumulates. The Gaussian treatment of the fluctuating phase gives a mean
residual amplitude $\exp(-\tfrac12\langle\Delta\phase^2\rangle)$ with
$\langle\Delta\phase^2\rangle\propto 2\tau/T_2$, yielding the stated
deficit; its monotonicity in $\tau$ and invariance under further pulses
follow from the monotonicity of the committed record.
\end{proof}

\subsection{The echo hierarchy and the engineering-independent floor}
\label{sec:hierarchy}

The refocusing idea has been pushed toward ever more complete reversal,
and the result is a hierarchy in which better engineering exposes a
sharper, engineering-independent floor.

\begin{itemize}[leftmargin=1.4em]
\item \textbf{Hahn echo} \citep{hahn1950}: reverses the static
single-spin inhomogeneity; floor at the homogeneous $T_2$, the deficit
$e^{-2\tau/T_2}$ of Corollary~\ref{cor:floor}.
\item \textbf{CPMG} \citep{carr1954,meiboom1958}: iterates the refocusing
pulse, converting slowly varying field components into refocusable
(configurational) dephasing before they commit; the decay time lengthens
toward, but never beyond, the limit set by fast fluctuations. Packing pulses
more densely does not breach the floor.
\item \textbf{Magic / polarization echo} \citep{rhim1971,levstein1998}:
engineers the many-body dipolar evolution to run as though under $-H$
rather than $+H$, the closest laboratory enactment of Loschmidt's
operation at the level of the interactions. The measured fidelity decay
becomes, in the strongly coupled regime, \emph{independent of the residual
imperfection} of the reversal: making the engineering more perfect does not
slow the decay \citep{levstein1998,jalabert2001,peres1984,gorin2006}.
\end{itemize}

\begin{remark}[The verdict]
\label{rem:verdict}
The reversibility presupposition predicts that the barrier to reversal is
contingent---improbability, imperfection, openness---and so should recede as
control improves. The echo hierarchy falsifies this prediction: the barrier
does not recede with control; it is intrinsic, set by the system's own
committed dynamics. This is precisely Theorem~\ref{thm:incoherence} and
Corollary~\ref{cor:floor} made experimental. The configurational part
refocuses arbitrarily well (Proposition~\ref{prop:echorecovery}); the
committed part does not refocus at all (Proposition~\ref{prop:refocus}); and
the limit of perfect reversal engineering is not perfect reversal but the
bare exposure of the floor. The experiment named after Loschmidt is the
experiment that decides against Loschmidt's premise.
\end{remark}

% =====================================================================
\section{Objections and Replies}
\label{sec:objections}
% =====================================================================

\paragraph{(O1) ``Erase the record, and the count is gone.''}
If one erases the visited-set record---or the observer's memory of it---the
states $S_0$ (before evolution) and $S_0'$ (after reversal) might seem
indistinguishable, hence the same. The reply is twofold. First, erasure is
itself an operation of positive duration and so increments the count
(Theorem~\ref{thm:incoherence}); it cannot decrement what it is an instance
of. Second, erasure is physically a commitment of information to the
environment, the irreversible kind of step of
Proposition~\ref{prop:refocus} and Corollary~\ref{cor:floor}: it enlarges,
not shrinks, the visited set of the joint system--environment partition.
Erasure removes the ability to \emph{read} the distinction from the system
alone; it does not remove the distinction
\citep{landauer1961,bennett1982,szilard1929}.

\paragraph{(O2) ``The count depends on the choice of section $\Sigma$.''}
Different sections yield different counts, but the monotonicity of each
(Theorem~\ref{thm:monotone}) and the incoherence of decrement
(Theorem~\ref{thm:incoherence}) hold for every positive-measure section.
The arrow is a property shared by all admissible counters, not an artefact
of one; this is the counterpart of the fact that all physical clocks agree
on the direction of time though not on its rate.

\paragraph{(O3) ``Virtual sub-states are a mere bookkeeping device.''}
They are a bookkeeping device, and they are physical. The off-shell
internal line is indispensable to the computed amplitude yet unobservable in
isolation (Theorem~\ref{thm:qftopacity}); the spectrometric off-shell
fraction is measured (Remark~\ref{rem:measuredfraction}); the
tailored-waveform cancellation is engineered (Remark~\ref{rem:swift}). A
device with measured fractions and engineered uses is not ``mere.'' What is
true is that virtual sub-states are confined to the reversible
configurational interior and are silent about the count---which is exactly
why they cannot be marshalled to reverse the arrow
(Theorem~\ref{thm:asymmetry}).

\paragraph{(O4) ``This is just coarse-graining renamed.''}
It is not. Coarse-graining \citep{gibbs1902} makes irreversibility relative
to a choice of smearing and leaves the fine-grained entropy constant by
Liouville's theorem. The present account makes no smearing: the visited set
$\Vis$ is defined by which cells the trajectory has \emph{entered}, a fact
about the actual path, not about a chosen resolution of the present state.
Two trajectories at the same present cell with different histories have
different visited sets; coarse-graining of the present state cannot tell
them apart, while $\Vis$ does (Remark~\ref{rem:negation}). The irreversibility
here is of the record, not of the resolution.

\paragraph{(O5) ``Why is $c$, the mass shell, or $[0,1]^m$ the right
admissible region?''}
The resolution does not depend on which admissible region one fixes; it
depends only on there \emph{being} one, so that ``off-shell'' is well
defined and the asymmetry of Section~\ref{sec:asymmetry} has content. The
three physical sections supply three concrete admissible regions---the mass
shell, the unit cube of normalised ion parameters, and the determined-phase
subspace of the spin ensemble---and the same abstract theorems specialise to
each.

% =====================================================================
\section{Conclusion}
\label{sec:conclusion}
% =====================================================================

Loschmidt's paradox is the appearance of contradiction between the
time-reversal symmetry of the microscopic equations and the manifest
asymmetry of macroscopic processes. We have argued that the appearance is
produced by collapsing two strata of a system's description that obey
different logics, and that holding them apart dissolves it without any
appeal to improbability, coarse-graining, or special boundary conditions.

The configurational stratum---the instantaneous microstate together with
the space of its decompositions---is genuinely reversible
(Theorem~\ref{thm:configreversible}), and it is rich, because it contains
virtual sub-states: decompositions whose components may be off-shell while
their composite is on-shell (Theorem~\ref{thm:virtualexist},
Theorem~\ref{thm:miracle}). The enumerative stratum---the count of
completed dynamical events---cannot decrease, because any operation that
would decrement it is itself a counted process
(Theorem~\ref{thm:incoherence}); it carries a visited-set entropy that is
monotone by the combinatorics of an appended record, with the second law as
a corollary (Corollary~\ref{cor:secondlaw}) and Boltzmann's $H$-theorem as
its statistical refinement. The two strata are logically independent
(Theorem~\ref{thm:distinct}), and the freedom of the first cannot reach the
second, because virtual sub-states are available only for the backward
completion of a fixed endpoint and never for forward construction
(Theorem~\ref{thm:asymmetry})---a structural asymmetry that is the
configurational shadow of the enumerative arrow. Velocity reversal acts on
the configurational stratum and leaves the enumerative one untouched, so it
restores configuration without restoring the past
(Theorem~\ref{thm:resolution}).

That virtual sub-states are not a formal contrivance is established by two
independent physical realisations and one controlled experiment. In
quantum field theory the off-shell internal line is a weighted, path-opaque
element of the reversible interior, and the unobservability of single
diagrams is the unobservability of single decompositions
(Theorem~\ref{thm:qftopacity}). In Fourier mass spectrometry a measured
fraction of an ion's natural decomposition lies off the physical manifold
while its mean recovers the physical coordinate
(Remark~\ref{rem:measuredfraction}). And the nuclear spin echo, the closest
laboratory approximation to Loschmidt's own operation, refocuses the
configurational part of the dynamics arbitrarily well while exposing an
engineering-independent floor exactly where information has been committed
(Corollary~\ref{cor:floor}, Remark~\ref{rem:verdict})---a seventy-year
experimental verdict against the assumption that time reversal is a coherent
operation on anything but configuration.

The arrow of time, in this account, is neither imposed by a boundary
condition nor purchased with a probability: it is the monotonicity of
counting, and the reason the richest reversible structure we possess cannot
overturn it is that this structure lives entirely within a single count and
is constitutively silent about which count one is at.

\section*{Acknowledgements}
The author thanks colleagues at the Chair of Proteomics and Bioanalytics
for discussions on Fourier-transform mass spectrometry.

\bibliographystyle{plainnat}
\bibliography{references}

\end{document}
